\section{引言}

数值逼近是一门数学与计算机的交叉学科。在计算机中，小数用浮点型数据存储，因此数值计算本质上都是近似计算。例如计算积分
\[
\int_{0}^{1} e^{-x^2}\,dx,
\]
一方面无法求出被积函数 $f(x)=e^{-x^2}$ 的原函数，另一方面积分值也只能用小数近似表示。这正体现了数值逼近的重要性。

此外，在求解多项式根时也需要近似方法。例如
\[
x^2+2x-1=0
\]
当然可用求根公式得到根式解，但仍需转化为近似小数；而五次以上方程没有一般的根式公式，更需发展系统的近似求根理论。

\subsection{数值逼近的一些简单应用}

数学分析中已经有许多数值逼近的例子：Taylor 展开、级数、函数项级数、Weierstrass 逼近定理、Fourier 级数等等。下面简单回顾几个分析中的特殊函数逼近。

\begin{example}
我们知道
\[
e = \sum_{k=0}^{\infty} \frac{1}{k!},
\]
这是一个无穷级数，实际计算只能取有限项。令 $a_n = \sum_{k=0}^{n} \frac{1}{k!}$，则 $\lim_{n\to\infty} a_n = e$。
研究 $a_n$ 趋近 $e$ 的速度，即分析误差项：
\[
e - a_n = \sum_{k=n+1}^{\infty} \frac{1}{k!}
      \le \frac{1}{(n+1)!}\left[1 + \frac{1}{n+1} + \frac{1}{(n+1)^2}+\dots\right]
      = \frac{1}{n\cdot n!}.
\]
这就是 $a_n$ 逼近 $e$ 的误差上界估计。
\end{example}

\begin{example}
证明 $\pi$ 的经典表达式
\[
\pi = 4\sum_{k=0}^{\infty}(-1)^k\frac{1}{2k+1}.
\]
由
\[
\frac{1}{1+x^2} = \sum_{k=0}^{\infty} (-1)^k x^{2k},
\]
且级数一致收敛，故
\[
\pi = 4\int_{0}^{1} \frac{1}{1+x^2}\,dx
    = 4\int_{0}^{1} \sum_{k=0}^{\infty} (-1)^k x^{2k}\,dx
    = 4\sum_{k=0}^{\infty} (-1)^k\frac{1}{2k+1}.
\]
\end{example}

\begin{example}
回到最初的积分
\[
\int_{0}^{1} e^{-x^2}\,dx,
\]
利用 $e^x = \sum_{k=0}^{\infty} \frac{x^k}{k!}$ 得
\[
\int_{0}^{1} e^{-x^2}\,dx = \sum_{k=0}^{\infty} \frac{(-1)^k}{k!} \int_{0}^{1} x^{2k}\,dx
                         = \sum_{k=0}^{\infty} \frac{(-1)^k}{(2k+1)k!},
\]
同样可得近似解，此处不再分析收敛速度。
\end{example}

通过以上例子可见 Taylor 展开与多项式级数在数值逼近中的重要性。我们先回顾 Taylor 展开的具体形式。

\subsection{泰勒展开在数值逼近中的应用}
设 $f$ 是 $\R$ 上无穷次可微的连续函数，则
\[
f(x) = \sum_{k=0}^{\infty} \frac{(x-x_0)^k}{k!} f^{(k)}(x_0)
\]
称为 $f$ 在 $x_0$ 处的 Taylor 展开式。Taylor 展开不一定收敛到原函数，但对初等函数如 $e^x$、$\ln x$、$\sin x$ 等均成立。一些常见函数在 $0$ 处的 Taylor 展开式如下：
\begin{enumerate}
    \item 指数函数 $f(x) = e^x$
    \[ e^x = \sum_{n=0}^{\infty}\frac{x^n}{n!}; \]
    \item 三角函数 $f(x) = \sin x$ 和 $g(x) = \cos x$
    \[
    \sin x =\sum_{n=0}^{\infty} \frac{(-1)^n}{(2n+1)!}x^{2n+1},\quad
    \cos x = \sum_{n=0}^{\infty} \frac{(-1)^n}{(2n)!}x^{2n};
    \]
    \item 对数函数 $f(x) = \ln (1+x)$
    \[ \ln (1+x) = \sum_{n=1}^{\infty} \frac{(-1)^{n-1}}{n} x^n; \]
    \item 反正切函数 $f(x) = \arctan x$
    \[ \arctan x = \sum_{n=0}^{\infty} \frac{(-1)^n}{2n+1} x^{2n+1}; \]
    \item 二项式函数 $f(x) = (1+x)^\alpha$
    \[ (1+x)^\alpha = \sum_{n=0}^{\infty} \frac{\alpha(\alpha-1)\cdots(\alpha-n+1)}{n!} x^n,\quad |x|<1. \]
\end{enumerate}

例如可通过代入 $x=2,\alpha=\frac12$ 求 $\sqrt{3}$ 的近似值。但计算机实际并不采用这种算法，因为收敛速度不够快。

\subsection{牛顿迭代法}
若从求方程根的角度，$\sqrt{3}$ 可视为 $f(x)=x^2-3$ 的零点。牛顿迭代法能提供更快的收敛。

\begin{theorem}
设 $f$ 在 $x^*$ 的邻域 $(x^*-\delta, x^*+\delta)$ 内二阶可导，$x_0$ 在该邻域内且与 $x^*$ 不重合，满足
\begin{enumerate}
    \item $x^*$ 是 $f$ 在该邻域内的唯一零点；
    \item $f'(x)$ 在 $x^*$ 的邻域内不为 $0$；
    \item $f''(x)$ 在 $x^*$ 的邻域内有界且不为 $0$。
\end{enumerate}
则递推式
\[ x_{n+1} = x_n - \frac{f(x_n)}{f'(x_n)} \]
产生的数列 $\{x_n\}$ 收敛到 $x^*$，且收敛速度是平方收敛。
\end{theorem}

\begin{proof}
令 $\Phi(x) = x - \frac{f(x)}{f'(x)}$，则 $x_{n+1} = \Phi(x_n)$。
由于
\[
\Phi'(x) = \frac{f(x)f''(x)}{[f'(x)]^2},
\]
在 $x^*$ 的邻域内只有 $x=x^*$ 使导数为零。故存在 $\delta'<\frac{\delta}{3}$，使得对任意 $x\in(x^*-\delta',x^*+\delta')$ 有 $|\Phi'(x)|\le \frac12$。
因此
\[
|x_{n+1}-x_n| = |\Phi(x_n)-\Phi(x_{n-1})| \le \frac12 |x_n-x_{n-1}|,
\]
从而 $\{x_n\}$ 收敛，且易证 $\{x_n\}\subset(x^*-\delta,x^*+\delta)$。
设 $x_n\to A$，由
\[
x_{n+1} = \Phi(x_n) = x_n - \frac{f(x_n)}{f'(x_n)} \to A - \frac{f(A)}{f'(A)} = A,
\]
得 $f(A)=0$，故 $A=x^*$。

令 $\varepsilon_n = x_n - x^*$，则
\[
x_{n+1}-x_n = -\frac{f(x_n)}{f'(x_n)},
\quad
\varepsilon_{n+1} - \varepsilon_n = -\frac{f(x_n)}{f'(x_n)},
\]
即
\[
\varepsilon_{n+1} = \frac{\varepsilon_n f'(x_n) - f(x_n)}{f'(x_n)}.
\]
利用 $f$ 在 $x^*$ 处的 Taylor 公式
\[
0 = f(x^*) = f(x_n) + f'(x_n)(-\varepsilon_n) + \frac{f''(\xi_n)}{2}\varepsilon_n^2,
\]
代入得
\[
\varepsilon_{n+1} = \frac{f''(\xi_n)}{2f'(x_n)}\varepsilon_n^2,
\]
可见误差平方收敛：存在常数 $C$ 使 $\left|\frac{\varepsilon_{n+1}}{\varepsilon_n^2}\right|\le C$。
\end{proof}

\begin{example}
用牛顿迭代法求 $\sqrt{3}$ 的近似值。取 $f(x)=x^2-3$，初值 $x_0=2$ 即可。
\end{example}

